\documentstyle[%


righttag%


,11pt%


,amssymb%


,russian%


,izvru%


]{amsart}





\newcommand{\tl}[3]{\medskip\noindent{\bf#1}\ #2 \dotfill #3 \par} 


\pagestyle{empty}


\begin{document}


%\tableofcontents


\par


\vskip 2cm


\par


\bigskip





\bigskip


\bigskip


\par


\centerline{\Large СОДЕРЖАНИЕ}


\bigskip


\bigskip


%%%%%%%%%% Use \newline %%%%%%%%%%%











\tl{Абдуллаев А.Р., Брагина Н.А.} 


{Операторы Грина с минимальной нормой}{3}


\tl{Власов В.В., Сакбаев В.Ж.}


{Корректная разрешимость дифференциальных уравнений\newline с последействием


в шкале пространств Соболева}{8}


\tl{Воскресенский Е.В.}


{Об аттракторах обыкновенных дифференциальных уравнений}{17}


\tl{Долгий Ю.Ф., Тарасян В.С.}


{Условия конечномерности оператора монодромии


для\newline периодических систем с последействием}{27}


\tl{Иванов А.Г.}


{К вопросу об оптимальном управлении 


почти периодическими движениями}{40}


\tl{Смолин Ю.Н.} 


{К вопросу об экспоненциальной устойчивости  


почти периодических\newline дифференциальных уравнений}{57}


\tl{Фишман Л.З.} 


{Определение опасных и безопасных границ  


области устойчивости\newline состояния равновесия  


уравнения второго порядка с запаздыванием}{61}


\tl{Чудинов К.М.} 


{Критерий устойчивости по части переменных 


автономной системы\newline дифференциальных уравнений}{67}


\tl{Щеглова А.А.}


{Левый регуляризирующий оператор для алгебро-дифференциальной\newline системы 


с запаздыванием}{73}














\vfill


\noindent\raisebox{2pt}{\copyright} Казанский госудаpственный унивеpситет


\end{document}


